\documentclass[10pt,a4paper]{amsart}
\usepackage{amsmath,amssymb,amsthm,graphicx}
\usepackage{verbatim}
\newenvironment{reference}{\comment}{\endcomment}
%\newenvironment{reference}{}{}
\newenvironment{slogan}{\comment}{\endcomment}
\newenvironment{history}{\comment}{\endcomment}
% For commutative diagrams we use Xy-pic
\usepackage[all]{xy}

% We use 2cell for 2-commutative diagrams.
\xyoption{2cell}
% Theorem environments.
%
\theoremstyle{plain}
\newtheorem{theorem}[subsection]{Theorem}
\newtheorem{proposition}[subsection]{Proposition}
\newtheorem{lemma}[subsection]{Lemma}

\theoremstyle{definition}
\newtheorem{definition}[subsection]{Definition}
\newtheorem{example}[subsection]{Example}
\newtheorem{exercise}[subsection]{Exercise}
\newtheorem{situation}[subsection]{Situation}

\theoremstyle{remark}
\newtheorem{remark}[subsection]{Remark}
\newtheorem{remarks}[subsection]{Remarks}

\numberwithin{equation}{subsection}

% Macros
%
\def\lim{\mathop{\mathrm{lim}}\nolimits}
\def\colim{\mathop{\mathrm{colim}}\nolimits}
\def\Spec{\mathop{\mathrm{Spec}}}
\def\Hom{\mathop{\mathrm{Hom}}\nolimits}
\def\Ext{\mathop{\mathrm{Ext}}\nolimits}
\def\SheafHom{\mathop{\mathcal{H}\!\mathit{om}}\nolimits}
\def\SheafExt{\mathop{\mathcal{E}\!\mathit{xt}}\nolimits}
\def\Sch{\mathit{Sch}}
\def\Mor{\mathop{\mathrm{Mor}}\nolimits}
\def\Ob{\mathop{\mathrm{Ob}}\nolimits}
\def\Sh{\mathop{\mathit{Sh}}\nolimits}
\def\NL{\mathop{N\!L}\nolimits}
\def\CH{\mathop{\mathrm{CH}}\nolimits}
\def\proetale{{pro\text{-}\acute{e}tale}}
\def\etale{{\acute{e}tale}}
\def\QCoh{\mathit{QCoh}}
\def\Ker{\mathop{\mathrm{Ker}}}
\def\Im{\mathop{\mathrm{Im}}}
\def\Coker{\mathop{\mathrm{Coker}}}
\def\Coim{\mathop{\mathrm{Coim}}}

% Boxtimes
%
\DeclareMathSymbol{\boxtimes}{\mathbin}{AMSa}{"02}

%
% Macros for moduli stacks/spaces
%
\def\QCohstack{\mathcal{QC}\!\mathit{oh}}
\def\Cohstack{\mathcal{C}\!\mathit{oh}}
\def\Spacesstack{\mathcal{S}\!\mathit{paces}}
\def\Quotfunctor{\mathrm{Quot}}
\def\Hilbfunctor{\mathrm{Hilb}}
\def\Curvesstack{\mathcal{C}\!\mathit{urves}}
\def\Polarizedstack{\mathcal{P}\!\mathit{olarized}}
\def\Complexesstack{\mathcal{C}\!\mathit{omplexes}}
% \Pic is the operator that assigns to X its picard group, usage \Pic(X)
% \Picardstack_{X/B} denotes the Picard stack of X over B
% \Picardfunctor_{X/B} denotes the Picard functor of X over B
\def\Pic{\mathop{\mathrm{Pic}}\nolimits}
\def\Picardstack{\mathcal{P}\!\mathit{ic}}
\def\Picardfunctor{\mathrm{Pic}}
\def\Deformationcategory{\mathcal{D}\!\mathit{ef}}
\usepackage{cleveref}
\crefname{originallemma}{original lemma}{original lemmas}
\Crefname{originallemma}{Original lemma}{Original lemmas}
\makeatletter
\newlabel{lemma-intersection-closed-in-quasi-compact}{{lemma-intersection-closed-in-quasi-compact}{source}{}{sourcepointer-1}{}}
\newlabel{lemma-intersection-closed-in-quasi-compact@cref}{{[originallemma][0][]lemma-intersection-closed-in-quasi-compact}{[1][0][]source}{}{}{}}
\makeatother
\crefname{originallemma}{original lemma}{original lemmas}
\Crefname{originallemma}{Original lemma}{Original lemmas}
\makeatletter
\newlabel{lemma-quasi-compact-in-Hausdorff}{{lemma-quasi-compact-in-Hausdorff}{source}{}{sourcepointer-2}{}}
\newlabel{lemma-quasi-compact-in-Hausdorff@cref}{{[originallemma][0][]lemma-quasi-compact-in-Hausdorff}{[1][0][]source}{}{}{}}
\makeatother
\title{The Baire category theorem for locally quasi-compact Hausdorff spaces}
\author{The Stacks Project Authors}
\date{stacks-project:topology.tex:local-snapshot-20261012; GNU FDL 1.2 or later}
\begin{document}
\maketitle
\noindent\textit{Editorial note.} All delivered mathematical text is the author's own, copied from the pinned Stacks Project chapter source sources/stacks-topology/pinned-original-source.tex; no statement or proof step was rewritten, reordered or removed.\par\medskip
\noindent\textit{Editorial note.} Preamble provenance: the Stacks Project keeps one single shared preamble for every chapter; this card takes it directly from the staged sources/stacks-topology/preamble.tex as three verbatim spans (lines 8-12, 14-18 and 154-217), so no derived file was written and nothing was added or redefined. The wrapper class and the base packages are supplied by the assembler.\par\medskip
\noindent\textit{Editorial note.} Bibliography provenance: neither the counted statement nor its proof contains a citation (this lane rescans the whole span before writing the card), so the card carries no bibliography and no entry was reconstructed.\par\medskip
\noindent\textit{Editorial note.} Selection provenance: this card was opened under the file-based control contract (CONTROL.md), with the mechanical pre-card gate run on its statement and proof interval before the card was written.\par\medskip
\noindent\textit{Allowed standard prerequisites (editorial).}\begin{itemize}
\item The statement that a quasi-compact subset of a Hausdorff space is closed. Invoked inside the counted proof (source lines 2282-2283). Makes the descending quasi-compact neighbourhoods closed.
\item The statement that an intersection of nonempty closed subsets of a quasi-compact space is nonempty. Invoked inside the counted proof (source lines 2284-2286). Produces a point in the total intersection.
\end{itemize}\medskip
\ifdefined\hypertarget\hypertarget{sourcepointer-1}{}\fi
\noindent\textit{Source reference.} Original lemma lemma-intersection-closed-in-quasi-compact: sources/stacks-topology/pinned-original-source.tex, source lines 1941--1941. Navigation pointer only: the counted text refers to the Stacks item lemma-intersection-closed-in-quasi-compact, which is an existing labeled result of this same chapter and is not part of this card. The pointer carries the target's own label as its number and its definition line as its span; its content is not reproduced here.\par\medskip
\ifdefined\hypertarget\hypertarget{sourcepointer-2}{}\fi
\noindent\textit{Source reference.} Original lemma lemma-quasi-compact-in-Hausdorff: sources/stacks-topology/pinned-original-source.tex, source lines 1894--1894. Navigation pointer only: the counted text refers to the Stacks item lemma-quasi-compact-in-Hausdorff, which is an existing labeled result of this same chapter and is not part of this card. The pointer carries the target's own label as its number and its definition line as its span; its content is not reproduced here.\par\medskip
\begin{lemma}[Baire category theorem]
\label{lemma-baire-category-locally-compact}
Let $X$ be a locally quasi-compact Hausdorff space.
Let $U_n \subset X$, $n \geq 1$ be dense open subsets. Then
$\bigcap_{n \geq 1} U_n$ is dense in $X$.
\end{lemma}

\begin{proof}
After replacing $U_n$ by $\bigcap_{i = 1, \ldots, n} U_i$
we may assume that $U_1 \supset U_2 \supset \ldots$.
Let $x \in X$. We will show that $x$ is in the closure of
$\bigcap_{n \geq 1} U_n$. Thus let $E$ be a neighbourhood of $x$.
To show that $E \cap \bigcap_{n \geq 1} U_n$ is nonempty we
may replace $E$ by a smaller neighbourhood. After replacing
$E$ by a smaller neighbourhood, we may assume that $E$ is quasi-compact.

\medskip\noindent
Set $x_0 = x$ and $E_0 = E$. Below, we will inductively choose
a point $x_i \in E_{i - 1} \cap U_i$ and a quasi-compact
neighbourhood $E_i$ of $x_i$ with $E_i \subset E_{i - 1} \cap U_i$.
Because $X$ is Hausdorff, the subsets $E_i \subset X$ are closed
(Lemma \ref{lemma-quasi-compact-in-Hausdorff}).
Since the $E_i$ are also nonempty we conclude that
$\bigcap_{i \geq 1} E_i$ is nonempty
(Lemma \ref{lemma-intersection-closed-in-quasi-compact}).
Since $\bigcap_{i \geq 1} E_i \subset E \cap \bigcap_{n \geq 1} U_n$
this proves the lemma.

\medskip\noindent
The base case $i = 0$ we have done above. Induction step.
Since $E_{i - 1}$ is a neighbourhood of $x_{i - 1}$ we can
find an open $x_{i - 1} \in W \subset E_{i - 1}$.
Since $U_i$ is dense in $X$
we see that $W \cap U_i$ is nonempty.
Pick any $x_i \in W \cap U_i$.
By definition of locally quasi-compact spaces we can
find a quasi-compact neighbourhood $E_i$ of $x_i$
contained in $W \cap U_i$. Then $E_i \subset E_{i - 1} \cap U_i$
as desired.
\end{proof}

\end{document}
